\clearpage
\renewcommand{\BookRunningTitle}{Inhalt.}
\SourceLeaf{537}\pdfbookmark[0]{Inhalt.}{historical-contents}
\begin{center}\Large INHALT.\par\medskip
\large GAUSS’ WERKE, BAND II. HÖHERE ARITHMETIK.
\end{center}
\BookRule{.25\textwidth}
\par\medskip\Needspace{5\baselineskip}\textit{Abhandlungen.}
\begingroup\fontsize{10}{13}\selectfont
\begin{longtable}{@{}>{\raggedright\arraybackslash}p{.60\linewidth}@{\quad}>{\raggedright\arraybackslash}p{.24\linewidth}@{\quad}r@{}}
Ein neuer Beweis eines arithmetischen Satzes .\,.\,.\,.\,.\,.\,.\,.\,.\,.\,.\,. & 1808 Jan. .\,. & Seite 1 \\[3pt]
Summation einiger besonderer Reihen .\,.\,.\,.\,.\,.\,.\,.\,.\,.\,. & 1808 Aug. .\,. & — 9 \\[3pt]
Neue Beweise und Erweiterungen des Fundamentaltheorems in der Theorie der quadratischen Reste .\,.\,.\,.\,.\,.\,.\,.\,.\,.\,.\,.\,.\,.\,.\,.\,. & 1817 Febr. . & — 47 \\[3pt]
Theorie der biquadratischen Reste. Erste Abhandlung .\kern.1em.\kern.4em.\,.\,.\,.\,.\,. & 1825 Apr. . & — 65 \\[3pt]
Theorie der biquadratischen Reste. Zweite Abhandlung .\,.\,.\,.\,.\,. & 1831 Apr. . & — 93 \\[3pt]
\end{longtable}\endgroup
\par\medskip\Needspace{5\baselineskip}\textit{Besprechungen eigener Schriften.}
\begingroup\fontsize{10}{13}\selectfont
\begin{longtable}{@{}>{\raggedright\arraybackslash}p{.60\linewidth}@{\quad}>{\raggedright\arraybackslash}p{.24\linewidth}@{\quad}r@{}}
Ein neuer Beweis eines arithmetischen Satzes .\,.\,.\,.\,.\,.\,.\,.\,.\,.\,.\,. & 1808 Mai .\,. & — 151 \\[3pt]
Summation einiger besonderer Reihen .\,.\,.\,.\,.\,.\,.\,.\,.\,.\,. & 1808 Sept. . & — 155 \\[3pt]
Fundamentaltheorem in der Theorie der Reste usw. .\,.\,.\,.\,.\,.\,.\,. & 1817 März . & — 159 \\[3pt]
Theorie der biquadratischen Reste. Abh. I. .\,.\,.\,.\,.\,.\,.\,.\,.\,. & 1825 April . & — 165 \\[3pt]
Theorie der biquadratischen Reste. Abh. II. .\,.\,.\,.\,.\,.\,.\,.\,.\,. & 1831 April . & — 169 \\[3pt]
\end{longtable}\endgroup
\par\medskip\Needspace{5\baselineskip}\textit{Besprechungen fremder Schriften.}
\begingroup\fontsize{10}{13}\selectfont
\begin{longtable}{@{}>{\raggedright\arraybackslash}p{.60\linewidth}@{\quad}>{\raggedright\arraybackslash}p{.24\linewidth}@{\quad}r@{}}
[\textsc{Dalberg}] Untersuchungen über die arithmetische und geometrische Irreduzibilität der Zahlen und ihrer Potenzen .\,.\,.\,.\,.\,.\,.\,.\,.\,.\,.\,.\,.\,.\,.\,. & 1809 März . & — 181 \\[3pt]
\textsc{Chernac}. Arithmetisches Sieb .\,.\,.\,.\,.\,.\,.\,.\,.\,.\,.\,.\,.\,.\,. & 1812 März . & — 181 \\[3pt]
\textsc{Burckhardt}. Teilertafeln .\,.\,.\,.\,.\,.\,. & 1814 Nov. 1816 Nov. 1817 Aug. . & — 183 \\[3pt]
\textsc{Erchinger}. Konstruktion des Siebzehnecks .\,.\,.\,.\,.\,.\,.\,.\,.\,.\,. & 1825 Dez. . & — 186 \\[3pt]
\textsc{Seeber}. Untersuchungen über die Eigenschaften der positiven ternären quadratischen Formen .\,.\,.\,.\,.\,.\,.\,.\,.\,.\,.\,.\,.\,.\,.\,.\,.\,. & 1831 Juli .\,. & — 188 \\[3pt]
\end{longtable}\endgroup
